\documentclass{labreport}

\usepackage{microtype}
\usepackage{enumitem}
\usepackage{pifont}
\usepackage{tabularx}
\usepackage{booktabs}

% Headheight adjustment for fancyhdr
\setlength{\headheight}{28pt}
\addtolength{\topmargin}{-13pt}

% Colors
\definecolor{codebg}{RGB}{248, 249, 251}
\definecolor{codeborder}{RGB}{218, 222, 229}
\definecolor{linkcolor}{RGB}{20, 60, 120}
\definecolor{zshgreen}{RGB}{0, 160, 80}
\definecolor{zshred}{RGB}{200, 30, 30}

% Hyperref setup with subtle link colors
\hypersetup{
    colorlinks=true,
    linkcolor=linkcolor,
    filecolor=linkcolor,
    urlcolor=linkcolor,
    citecolor=linkcolor
}

% Code and terminal listings
\lstdefinestyle{bash}{
    basicstyle=\ttfamily\footnotesize,
    backgroundcolor=\color{codebg},
    frame=single,
    framerule=0.5pt,
    rulecolor=\color{codeborder},
    breaklines=true,
    breakatwhitespace=false,
    showstringspaces=false,
    keepspaces=true,
    columns=flexible,
    numbers=none,
    xleftmargin=3pt,
    xrightmargin=3pt,
    framesep=4pt
}

\lstdefinestyle{c}{
    language=C,
    basicstyle=\ttfamily\footnotesize,
    backgroundcolor=\color{codebg},
    frame=single,
    framerule=0.5pt,
    rulecolor=\color{codeborder},
    keywordstyle=\color{blue!85!black}\bfseries,
    commentstyle=\color{gray!80!black}\itshape,
    stringstyle=\color{teal!80!black},
    numbers=left,
    numberstyle=\tiny\color{gray!80!black},
    numbersep=8pt,
    breaklines=true,
    breakatwhitespace=false,
    showstringspaces=false,
    keepspaces=true,
    columns=flexible,
    tabsize=2,
    xleftmargin=14pt,
    xrightmargin=3pt,
    framesep=4pt
}

\lstdefinestyle{asm}{
    language=[x86masm]Assembler,
    basicstyle=\ttfamily\footnotesize,
    backgroundcolor=\color{codebg},
    frame=single,
    framerule=0.5pt,
    rulecolor=\color{codeborder},
    keywordstyle=\color{purple!85!black}\bfseries,
    commentstyle=\color{gray!80!black}\itshape,
    stringstyle=\color{teal!80!black},
    numbers=left,
    numberstyle=\tiny\color{gray!80!black},
    numbersep=8pt,
    breaklines=true,
    breakatwhitespace=false,
    showstringspaces=false,
    keepspaces=true,
    columns=flexible,
    xleftmargin=14pt,
    xrightmargin=3pt,
    framesep=4pt
}

\lstset{
    extendedchars=true,
    literate=
        {➜}{{\textcolor{zshgreen}{\ding{220}}}}{1}
        {├─}{|-\,}{2}
        {└─}{\`-\,}{2}
        {├}{|}{1}
        {─}{-}{1}
        {└}{\`}{1}
        {✗}{{\textcolor{zshred}{\ding{55}}}}{1}
}

% Section styling for clean document structure
\titleformat{\section}
  {\large\bfseries}{\thesection}{0.8em}{}[\vspace{1pt}\hrule height 0.5pt]
\titleformat{\subsection}
  {\normalsize\bfseries}{\thesubsection}{0.8em}{}
\titleformat{\subsubsection}
  {\small\bfseries}{\thesubsubsection}{0.8em}{}

\begin{document}

\labtitle
    {18}
    {A User-Level Thread Switch}
    {st. gr. FAF-243}
    {Șoimu Ionuț}
    {asist. univ.\\[0.15cm] Reițman Patricia}

\tableofcontents
\newpage

% ==============================================================================
\section{Lab 18: Track B -- A User-Level Thread Switch (xv6)}

\subsection{Before you start}

\subsubsection*{Purpose and Scope}
In operating systems architecture, thread management can occur at two distinct abstraction layers:
\begin{enumerate}[label=(\alph*), noitemsep]
    \item \textbf{Kernel-Level Threads (1:1 Model):} The kernel scheduler creates, tracks, and switches execution contexts between threads using hardware timer interrupts and supervisor-mode state frames (\texttt{struct trapframe}).
    \item \textbf{User-Level Threads (M:1 Model / Green Threads / Coroutines):} Thread scheduling, context persistence, and execution switching occur entirely within user space code, without issuing system calls or transitioning CPU privilege rings into supervisor mode.
\end{enumerate}

The objective of Laboratory Work 18 is to design and implement a compact, high-performance \textbf{user-level cooperative threading package} (\texttt{uthread}) in the \textbf{xv6-riscv} operating system. Specifically:
\begin{enumerate}[label=(\alph*), noitemsep]
    \item Allocate dedicated thread control structures and private execution call stacks for independent threads;
    \item Define the execution \texttt{context} reflecting the RISC-V calling convention;
    \item Implement thread initialization and bootstrapping in \texttt{thread\_create()};
    \item Implement cooperative scheduling and context handoff in \texttt{thread\_schedule()} and \texttt{thread\_yield()};
    \item Author the machine-level register preservation routine in assembly (\texttt{user/uthread\_switch.S});
    \item Integrate the build targets into the xv6 build system without modifying kernel-space code;
    \item Boot xv6 under QEMU and empirically verify three concurrent threads alternating execution until completion.
\end{enumerate}

% ==============================================================================
\subsection{Part 1. The switch}

\subsubsection{1.1 Thread Control Block and Context: user/uthread.c}
A thread requires two core resources to run: its own stack and saved CPU state when not currently executing. 

\paragraph{Callee-Saved Registers and Context Definition:}
Under the standard RISC-V calling convention (RV64 ABI), registers are partitioned into caller-saved and callee-saved:
\begin{itemize}[noitemsep]
    \item \textbf{Caller-saved registers} (\texttt{ra}, \texttt{t0}--\texttt{t6}, \texttt{a0}--\texttt{a7}): If a function expects these values to persist across a call, the compiler generates instructions to push them onto the calling thread's own stack prior to invoking \texttt{thread\_switch()}.
    \item \textbf{Callee-saved registers} (\texttt{ra}, \texttt{sp}, \texttt{s0}--\texttt{s11}): The called function (\texttt{thread\_switch}) is obligated to preserve these across execution.
\end{itemize}

Therefore, our \texttt{struct context} requires exactly $14$ 64-bit integer fields:

\begin{lstlisting}[style=c]
/* the registers a thread must keep across a switch (callee-saved) */
struct context {
  uint64 ra;
  uint64 sp;

  // callee-saved registers
  uint64 s0;
  uint64 s1;
  uint64 s2;
  uint64 s3;
  uint64 s4;
  uint64 s5;
  uint64 s6;
  uint64 s7;
  uint64 s8;
  uint64 s9;
  uint64 s10;
  uint64 s11;
};
\end{lstlisting}

\paragraph{Thread Representation:}
Each thread is encapsulated in a \texttt{struct thread}:
\begin{lstlisting}[style=c]
struct thread {
  struct context ctx;                                             /* saved registers */
  char           stack[STACK_SIZE] __attribute__((aligned(16)));  /* private stack */
  int            state;                                           /* FREE, RUNNING, RUNNABLE */
};
\end{lstlisting}
The stack must be strictly aligned to a 16-byte boundary as mandated by the RISC-V standard stack frame alignment rules.

\paragraph{Thread Creation: thread\_create()}
When a new thread is spawned, an unused entry in \texttt{all\_thread} is located and transitioned to the \texttt{RUNNABLE} state. To prepare the thread for its first execution switch, two crucial context fields must be initialized:
\begin{enumerate}[noitemsep]
    \item \texttt{ctx.ra}: Initialized to the memory address of the target function (\texttt{func}). When \texttt{thread\_switch} executes \texttt{ret}, the program counter will jump directly to \texttt{func}.
    \item \texttt{ctx.sp}: Initialized to the top of the thread's stack (\texttt{t->stack + STACK\_SIZE}). Because stacks on RISC-V grow downwards toward lower addresses, the initial stack pointer must reside at the highest valid address of the allocated buffer.
\end{enumerate}

\begin{lstlisting}[style=c]
void
thread_create(void (*func)())
{
  struct thread *t;

  for (t = all_thread; t < all_thread + MAX_THREAD; t++) {
    if (t->state == FREE) break;
  }
  t->state = RUNNABLE;
  t->ctx.ra = (uint64)func;
  t->ctx.sp = (uint64)(t->stack + STACK_SIZE);
}
\end{lstlisting}

\paragraph{Cooperative Scheduling: thread\_schedule() and thread\_yield()}
A round-robin scheduler inspects the thread table to locate the next \texttt{RUNNABLE} thread. When switching from \texttt{current\_thread} to \texttt{next\_thread}:
\begin{lstlisting}[style=c]
  if (current_thread != next_thread) {
    /* switch threads */
    next_thread->state = RUNNING;
    struct thread *old = current_thread;
    current_thread = next_thread;
    thread_switch(&old->ctx, &next_thread->ctx);
  }
\end{lstlisting}

Voluntary CPU relinquishment occurs when a thread invokes \texttt{thread\_yield()}, setting its state to \texttt{RUNNABLE} and entering the scheduler.

\paragraph{Complete Implementation of user/uthread.c:}
\begin{lstlisting}[style=c]
#include "kernel/types.h"
#include "user/user.h"

/* Possible states of a thread: */
#define FREE        0x0
#define RUNNING     0x1
#define RUNNABLE    0x2

#define STACK_SIZE  8192
#define MAX_THREAD  4

struct context {
  uint64 ra;
  uint64 sp;
  uint64 s0;
  uint64 s1;
  uint64 s2;
  uint64 s3;
  uint64 s4;
  uint64 s5;
  uint64 s6;
  uint64 s7;
  uint64 s8;
  uint64 s9;
  uint64 s10;
  uint64 s11;
};

struct thread {
  struct context ctx;
  char           stack[STACK_SIZE] __attribute__((aligned(16)));
  int            state;
};

struct thread all_thread[MAX_THREAD];
struct thread *current_thread;
extern void thread_switch(struct context *old, struct context *new);

void
thread_init(void)
{
  current_thread = &all_thread[0];
  current_thread->state = RUNNING;
}

void
thread_schedule(void)
{
  struct thread *t, *next_thread;

  next_thread = 0;
  t = current_thread + 1;
  for(int i = 0; i < MAX_THREAD; i++){
    if(t >= all_thread + MAX_THREAD)
      t = all_thread;
    if(t->state == RUNNABLE) {
      next_thread = t;
      break;
    }
    t = t + 1;
  }

  if (next_thread == 0) {
    printf("thread_schedule: no runnable threads\n");
    exit(-1);
  }

  if (current_thread != next_thread) {
    next_thread->state = RUNNING;
    struct thread *old = current_thread;
    current_thread = next_thread;
    thread_switch(&old->ctx, &next_thread->ctx);
  }
}

void
thread_create(void (*func)())
{
  struct thread *t;

  for (t = all_thread; t < all_thread + MAX_THREAD; t++) {
    if (t->state == FREE) break;
  }
  t->state = RUNNABLE;
  t->ctx.ra = (uint64)func;
  t->ctx.sp = (uint64)(t->stack + STACK_SIZE);
}

void
thread_yield(void)
{
  current_thread->state = RUNNABLE;
  thread_schedule();
}

volatile int a_started, b_started, c_started;
volatile int a_n, b_n, c_n;

void
thread_a(void)
{
  int i;
  printf("thread_a started\n");
  a_started = 1;
  while(b_started == 0 || c_started == 0)
    thread_yield();
  for (i = 0; i < 100; i++) {
    printf("thread_a %d\n", i);
    a_n += 1;
    thread_yield();
  }
  printf("thread_a: exit after %d\n", a_n);
  current_thread->state = FREE;
  thread_schedule();
}

void
thread_b(void)
{
  int i;
  printf("thread_b started\n");
  b_started = 1;
  while(a_started == 0 || c_started == 0)
    thread_yield();
  for (i = 0; i < 100; i++) {
    printf("thread_b %d\n", i);
    b_n += 1;
    thread_yield();
  }
  printf("thread_b: exit after %d\n", b_n);
  current_thread->state = FREE;
  thread_schedule();
}

void
thread_c(void)
{
  int i;
  printf("thread_c started\n");
  c_started = 1;
  while(a_started == 0 || b_started == 0)
    thread_yield();
  for (i = 0; i < 100; i++) {
    printf("thread_c %d\n", i);
    c_n += 1;
    thread_yield();
  }
  printf("thread_c: exit after %d\n", c_n);
  current_thread->state = FREE;
  thread_schedule();
}

int
main(int argc, char *argv[])
{
  a_started = b_started = c_started = 0;
  a_n = b_n = c_n = 0;
  thread_init();
  thread_create(thread_a);
  thread_create(thread_b);
  thread_create(thread_c);
  current_thread->state = FREE;
  thread_schedule();
  exit(0);
}
\end{lstlisting}

% ------------------------------------------------------------------------------
\subsubsection{1.2 Low-Level Context Switch: user/uthread\_switch.S}
The context switch operation cannot be expressed in standard C because saving and restoring arbitrary general-purpose registers would corrupt the compiler's own frame. We author \texttt{thread\_switch} directly in assembly.

The function takes two pointer arguments passed in registers \texttt{a0} and \texttt{a1}:
\begin{itemize}[noitemsep]
    \item \texttt{a0}: \texttt{struct context *old} (pointer to the outgoing thread's context structure);
    \item \texttt{a1}: \texttt{struct context *new} (pointer to the incoming thread's context structure).
\end{itemize}

It executes two symmetric phases:
\begin{enumerate}[noitemsep]
    \item \textbf{Save Phase (\texttt{sd}):} Writes all 14 callee-saved registers to offsets \texttt{0(a0)} through \texttt{104(a0)}.
    \item \textbf{Restore Phase (\texttt{ld}):} Reads the 14 callee-saved registers from offsets \texttt{0(a1)} through \texttt{104(a1)}.
    \item \textbf{Return (\texttt{ret}):} Executes \texttt{jalr zero, 0(ra)}, jumping into the address loaded into \texttt{ra}.
\end{enumerate}

\begin{lstlisting}[style=asm]
	.text

	/*
	 * save the old thread's registers,
	 * restore the new thread's registers.
	 */

	.globl thread_switch
thread_switch:
	sd ra, 0(a0)
	sd sp, 8(a0)
	sd s0, 16(a0)
	sd s1, 24(a0)
	sd s2, 32(a0)
	sd s3, 40(a0)
	sd s4, 48(a0)
	sd s5, 56(a0)
	sd s6, 64(a0)
	sd s7, 72(a0)
	sd s8, 80(a0)
	sd s9, 88(a0)
	sd s10, 96(a0)
	sd s11, 104(a0)

	ld ra, 0(a1)
	ld sp, 8(a1)
	ld s0, 16(a1)
	ld s1, 24(a1)
	ld s2, 32(a1)
	ld s3, 40(a1)
	ld s4, 48(a1)
	ld s5, 56(a1)
	ld s6, 64(a1)
	ld s7, 72(a1)
	ld s8, 80(a1)
	ld s9, 88(a1)
	ld s10, 96(a1)
	ld s11, 104(a1)

	ret    /* return to ra */
\end{lstlisting}

% ------------------------------------------------------------------------------
\subsubsection{1.3 Build Automation: Makefile Integration}
To link \texttt{uthread.o} and \texttt{uthread\_switch.o} into the final xv6 executable \texttt{\_uthread}, the following rules were incorporated into \texttt{Makefile}:

\begin{lstlisting}[style=bash]
$U/uthread_switch.o: $U/uthread_switch.S
	$(CC) $(CFLAGS) -c -o $U/uthread_switch.o $U/uthread_switch.S

$U/_uthread: $U/uthread.o $U/uthread_switch.o $(ULIB) $U/user.ld
	$(LD) $(LDFLAGS) -T $U/user.ld -o $U/_uthread $U/uthread.o $U/uthread_switch.o $(ULIB)
	$(OBJDUMP) -S $U/_uthread > $U/uthread.asm
\end{lstlisting}

And \texttt{\$U/\_uthread} was appended to the xv6 user utilities list (\texttt{UPROGS}).

% ==============================================================================
\subsection{Part 2. Run it and verify}

\subsubsection{Compilation and QEMU Execution}
The kernel and user filesystem image were rebuilt:
\begin{lstlisting}[style=bash]
$ make user/_uthread fs.img
riscv64-linux-gnu-gcc -Wall -Werror ... -c -o user/uthread.o user/uthread.c
riscv64-linux-gnu-gcc -Wall -Werror ... -c -o user/uthread_switch.o user/uthread_switch.S
riscv64-linux-gnu-ld -z max-page-size=4096 -T user/user.ld -o user/_uthread user/uthread.o user/uthread_switch.o user/ulib.o user/usys.o user/printf.o user/umalloc.o
riscv64-linux-gnu-objdump -S user/_uthread > user/uthread.asm
mkfs/mkfs fs.img README ... user/_uthread
\end{lstlisting}

We then booted xv6 inside QEMU and executed \texttt{uthread}. Below is the real terminal transcript:

\begin{lstlisting}[style=bash]
$ uthread
thread_a started
thread_b started
thread_c started
thread_c 0
thread_a 0
thread_b 0
thread_c 1
thread_a 1
thread_b 1
thread_c 2
thread_a 2
thread_b 2
thread_c 3
thread_a 3
thread_b 3
...
thread_c 98
thread_a 98
thread_b 98
thread_c 99
thread_a 99
thread_b 99
thread_c: exit after 100
thread_a: exit after 100
thread_b: exit after 100
thread_schedule: no runnable threads
$ 
\end{lstlisting}

All three threads synchronized startup using their \texttt{volatile} flags, alternated counting turns in round-robin sequence from 0 to 99, completed their loops at 100, and exited smoothly. When all threads completed, the scheduler properly reported \texttt{thread\_schedule: no runnable threads} and terminated.

% ==============================================================================
\subsection{Observe: Deep Dive Analysis}

\begin{quote}
\textbf{Observe:} \textit{What exactly does \texttt{thread\_switch} save and load? Why is \texttt{ra} the register that makes execution continue in the other thread?}
\end{quote}

\subsubsection{1. Exactly What Registers Are Saved and Loaded}
\texttt{thread\_switch} saves and restores precisely \textbf{14 registers} (112 bytes total) defined by the RISC-V calling convention as \textbf{callee-saved}:
\begin{center}
\begin{tabularx}{\textwidth}{l l l X}
\toprule
\textbf{Register} & \textbf{ABI Name} & \textbf{Width} & \textbf{Architectural Role} \\
\midrule
\texttt{x1}  & \texttt{ra}     & 64-bit & Return address (target instruction for \texttt{ret}) \\
\texttt{x2}  & \texttt{sp}     & 64-bit & Stack pointer (top of thread's private stack frame) \\
\texttt{x8}  & \texttt{s0/fp}  & 64-bit & Saved register 0 / frame pointer \\
\texttt{x9}  & \texttt{s1}     & 64-bit & Saved register 1 \\
\texttt{x18}--\texttt{x27} & \texttt{s2}--\texttt{s11} & 64-bit each & Saved registers 2 through 11 \\
\bottomrule
\end{tabularx}
\end{center}

\paragraph{Why Caller-Saved Registers Are Omitted:}
Notice that caller-saved registers (\texttt{a0}--\texttt{a7}, \texttt{t0}--\texttt{t6}) are \textbf{not} saved by \texttt{thread\_switch}. Under the C ABI, because \texttt{thread\_switch} is invoked as an ordinary C function call (\texttt{thread\_switch(\&old->ctx, \&new->ctx)}), the C compiler itself is responsible for pushing any active caller-saved variables onto the thread's stack before the call occurs. Thus, saving only the callee-saved set plus \texttt{sp} and \texttt{ra} is necessary and sufficient to restore full execution context.

\subsubsection{2. Why ra Makes Execution Continue in the Other Thread}
The \texttt{ra} (Return Address) register holds the 64-bit program address to which the CPU branches when executing a subroutine return (\texttt{ret}, which is pseudo-assembly for \texttt{jalr zero, 0(ra)}).

In our user threading engine, \texttt{ra} governs thread continuation in two distinct operational phases:

\begin{enumerate}[leftmargin=*]
    \item \textbf{First Time a Thread Runs (Bootstrapping via \texttt{thread\_create}):}\\
    When a thread is initialized in \texttt{thread\_create(func)}, its saved context is explicitly configured:
    \begin{center}
        \texttt{t->ctx.ra = (uint64)func;}
    \end{center}
    When the scheduler subsequently selects this thread for the first time, \texttt{thread\_switch} executes:
    \begin{center}
        \texttt{ld ra, 0(a1)} $\longrightarrow$ loads address of \texttt{func} into register \texttt{ra}\\
        \texttt{ld sp, 8(a1)} $\longrightarrow$ loads top of thread's stack into register \texttt{sp}\\
        \texttt{ret} $\longrightarrow$ CPU branches directly to \texttt{func}
    \end{center}
    Execution begins immediately inside \texttt{func} on the thread's dedicated stack.

    \item \textbf{Subsequent Context Switches (Resumption after \texttt{thread\_yield}):}\\
    When an already-running thread yields the CPU, it calls \texttt{thread\_schedule()}, which in turn issues:
    \begin{center}
        \texttt{thread\_switch(\&old->ctx, \&next\_thread->ctx);}
    \end{center}
    The assembly instruction \texttt{call thread\_switch} automatically stores the return address (the next instruction inside \texttt{thread\_schedule} immediately following the call) into register \texttt{ra}.
    \begin{itemize}[noitemsep]
        \item When \texttt{thread\_switch} starts, \texttt{sd ra, 0(a0)} saves this instruction pointer into the yielding thread's context structure.
        \item When this thread is later rescheduled, \texttt{ld ra, 0(a1)} restores this saved pointer back into register \texttt{ra}.
        \item When \texttt{thread\_switch} executes \texttt{ret}, the CPU returns to the exact location inside \texttt{thread\_schedule} where the thread paused, seamlessly resuming user code execution.
    \end{itemize}
\end{enumerate}

Thus, \texttt{ra} serves as the pivotal register that directs hardware instruction fetching to the exact point of suspension or initial entry.

% ==============================================================================
\subsection{Conclusion}
Laboratory Work 18 demonstrated that thread scheduling and context switching are not proprietary magical operations reserved solely for the operating system kernel. By preserving and restoring the callee-saved registers (\texttt{ra}, \texttt{sp}, \texttt{s0}--\texttt{s11}) in assembly, an application can switch execution threads entirely in user space in just a few dozen CPU cycles. This architecture forms the foundational design of modern async runtimes, goroutines, and lightweight fiber libraries.

\end{document}
